\section{Implicit relevance and appropriate influence}\label{sec:C27:influence}
InMind tests a failure mode in which a memory is relevant through a knowledge bridge rather than lexical resemblance~\cite{inmind2026}. Its paired controls distinguish storage, bridging knowledge, and query-conditioned retrieval. This is more diagnostic than asking only whether the agent can repeat a fact on demand. The benchmark is selective, however, and its authors identify domain coverage, shared answerer-judge identity, and missing negative controls against over-warning as limitations.

The mechanism suggests a concrete comparison. Hold the decisive memory and task fixed, then vary explicit recall queries, indirect application queries, direct in-context exposure, and retrieval-mediated exposure. A gap between the latter two implicates the access interface rather than necessarily the model's knowledge. To test whether a proposed remedy overuses memory, add cases where a similar stored fact should not affect the answer. Without those negatives, more frequent warning or personalization can look like better judgment.

MemCalib isolates another stage: the amount of influence a supplied memory should have~\cite{memcalib2026}. It distinguishes over-use from under-use and studies training signals that balance the two. Because the relevant memory blocks are supplied, the result concerns memory-use behavior rather than complete storage and retrieval quality. Its counterfactual localization is a credit proxy with its own uncertainty, not a direct measurement of every token's true causal contribution.

An influence evaluation needs context-dependent targets. A memory can be evidence, a preference, a constraint, or an obsolete observation. The proper response depends on provenance, current authority, and task scope. A simple rule to always prioritize memory is therefore inadequate. Test contradictory updates, exceptions, and cases where a retrieved instruction lacks authority to modify the user's objective.

These studies jointly motivate separating recall from decision use. Better retrieval can expose more relevant evidence while a poorly calibrated reader overweights it. Better memory-use training can improve interpretation while retrieval still misses implicit relevance. A complete experiment crosses retrieval policy with reader policy and reports their interaction. This factorial design can reveal whether an improvement transfers across memory engines or depends on a particular writer-reader pairing.

